% !TeX root = tfg.tex
% !TeX encoding = utf8
%
%*******************************************************
% Preámbulo propio: paquetes y macros añadidos a la clase ugrTFG
% (se carga justo después de \documentclass; no modificar ugrTFG.cls)
%*******************************************************

% -------------------------------------------------------------------
% Versión borrador / final
% -------------------------------------------------------------------
% Con \borradortrue se muestran las marcas [PENDIENTE] y [REVISAR];
% con \borradorfalse desaparecen en la versión final.
\newif\ifborrador
\borradortrue

\ifborrador
  \newcommand{\pendiente}[1]{\textcolor{red}{\textbf{[PENDIENTE: #1]}}}
  \newcommand{\revisar}[1]{\textcolor{orange}{\textbf{[REVISAR: #1]}}}
\else
  \newcommand{\pendiente}[1]{}
  \newcommand{\revisar}[1]{}
\fi

% -------------------------------------------------------------------
% Paquetes
% -------------------------------------------------------------------
\usepackage{siunitx}
\sisetup{output-decimal-marker={,}, group-separator={.}, group-minimum-digits=4}
\usepackage{subcaption}
\usepackage{multirow}
\usepackage{tabularx}
\usepackage{longtable}
\usepackage{array}
\usepackage{pgfgantt}

% cleveref siempre después de hyperref (que carga la clase)
\usepackage[spanish, noabbrev]{cleveref}
\crefname{teorema}{teorema}{teoremas}
\Crefname{teorema}{Teorema}{Teoremas}
\crefname{proposicion}{proposición}{proposiciones}
\Crefname{proposicion}{Proposición}{Proposiciones}
\crefname{lema}{lema}{lemas}
\Crefname{lema}{Lema}{Lemas}
\crefname{corolario}{corolario}{corolarios}
\Crefname{corolario}{Corolario}{Corolarios}
\crefname{definicion}{definición}{definiciones}
\Crefname{definicion}{Definición}{Definiciones}
\crefname{observacion}{observación}{observaciones}
\Crefname{observacion}{Observación}{Observaciones}
\crefname{ejemplo}{ejemplo}{ejemplos}
\Crefname{ejemplo}{Ejemplo}{Ejemplos}
\crefname{part}{parte}{partes}
\Crefname{part}{Parte}{Partes}

% -------------------------------------------------------------------
% Entornos
% -------------------------------------------------------------------
% La clase no define "demostracion": se construye sobre proof de amsthm,
% que con babel en español ya imprime «Demostración».
\newenvironment{demostracion}{\begin{proof}}{\end{proof}}

% -------------------------------------------------------------------
% Notación recurrente (definida una sola vez)
% -------------------------------------------------------------------
% Conjuntos y probabilidad
\newcommand{\R}{\mathbb{R}}
\newcommand{\N}{\mathbb{N}}
\DeclareMathOperator{\E}{\mathbb{E}}
\newcommand{\Prob}{\mathbb{P}}
\newcommand{\Ind}{\mathbb{1}}
\DeclareMathOperator*{\argmax}{arg\,max}
\DeclareMathOperator*{\argmin}{arg\,min}
\DeclareMathOperator{\softmax}{softmax}
\DeclareMathOperator{\Var}{Var}

% Distribuciones y conjuntos ID / OOD
\newcommand{\Pin}{P_{\mathrm{in}}}
\newcommand{\Pout}{P_{\mathrm{out}}}
\newcommand{\Din}{\mathcal{D}_{\mathrm{in}}}
\newcommand{\Dout}{\mathcal{D}_{\mathrm{out}}}
\newcommand{\Xin}{X_{\mathrm{in}}}
\newcommand{\Xout}{X_{\mathrm{out}}}

% Métricas (válidas en texto y en modo matemático)
\newcommand{\AUROC}{\ensuremath{\mathrm{AUROC}}\xspace}
\newcommand{\AUPRin}{\ensuremath{\mathrm{AUPR\text{-}In}}\xspace}
\newcommand{\AUPRout}{\ensuremath{\mathrm{AUPR\text{-}Out}}\xspace}
\newcommand{\FPR}{\ensuremath{\mathrm{FPR}}\xspace}
\newcommand{\TPR}{\ensuremath{\mathrm{TPR}}\xspace}
\newcommand{\CCR}{\ensuremath{\mathrm{CCR}}\xspace}
\newcommand{\FPRn}[1]{\ensuremath{\mathrm{FPR@}#1}\xspace}
\newcommand{\CCRn}[1]{\ensuremath{\mathrm{CCR@FPR}#1}\xspace}

% Nombres de métodos y conjuntos (formato consistente en el texto)
\newcommand{\stoodx}{STOOD-X\xspace}
\newcommand{\nctcrc}{NCT-CRC-HE-100K-NONORM\xspace}
\newcommand{\clase}[1]{\texttt{#1}}

\endinput
